\section{收盘行情与流动性}
2026年7月14日通鼎互联收盘20.86元，较前收19.87元上涨4.98\%，成交量1.2015亿股，成交额24.04亿元，日内高点20.86元、低点19.12元。成交额和换手处于事件交易状态：这是H1预告之后的高流动性博弈日，不是低位安静吸筹的证据。

\begin{exhibitbox}[表6-1\quad 资金层级]
\begin{tabularx}{\exhibitboxwidth}{L{2.7cm}X L{2.9cm}X}
\textbf{资金层} & \textbf{可确认事实} & \textbf{资金性质} & \textbf{结论} \\
控制股东 & 通鼎集团31.51\%、沈小平4.55\%（Q1末） & 稳定控制块 & 不是边际交易资金 \\
公募/外资托管 & 华商两只基金、易方达、HKSCC、高盛出现在Q1前十 & 季度末机构持仓 & 有机构基础，不代表7月仍持有 \\
机构专用席位 & 4/9、4/16、5/12、6/4、6/8、6/17和7/1均有事件记录，方向交替 & 主动交易机构 & 机构参与但分歧大 \\
融资资金 & 7/13融资余额约12.17亿元，净偿还约0.93亿元 & 杠杆资金 & 放大波动，不等于长期看多 \\
\end{tabularx}
\sourcenote{data/capital\_structure\_20260714.md；sources/market-20260714/lhb.html；2026Q1报告。}
\end{exhibitbox}

\section{机构票还是游资票}
它有真实机构底仓，但边际价格由机构、融资、量化和活跃席位共同推动。龙虎榜数据显示，6月4日、6月8日和6月17日机构净买入，7月1日连续三日下跌事件中机构净卖出约2.07亿元；这种买卖交替更像事件/趋势交易，而不是单边长期吸筹。机构专用席位不能穿透到具体基金，活跃营业部也不能直接等同游资身份。

因此分类为“机构参与的高换手趋势/事件票”，不是纯机构票，也不是纯游资票。资金评分为：机构底仓73/100、主动交易92/100、拥挤度88/100、流动性94/100、综合风格86/100。评分衡量证据结构和交易行为，不预测未来涨跌。

\section{支撑、阻力与交易纪律}
\begin{itemize}
\item CNY18.50：市场情绪锚，跌破代表H1事件溢价开始收敛。
\item CNY15--16：前期业绩重估区，只有基本面未破坏时才有观察意义。
\item CNY21.86：按当日价格规则计算的公开涨停参考位。
\item CNY25--30：前期事件交易供给区，不把突破本身视为基本面确认。
\end{itemize}
这些水平是行为参考而非确定技术信号；融资余额、H1现金流和正式半年报是更重要的风控变量。
